\section{问题重述}
\subsection{问题背景}
\begin{enumerate}[label=(\arabic*), leftmargin=2.5em, itemsep=0.3em]
\item \textbf{不要直接复制赛题。} 应提炼题目的背景、已知条件和研究任务，用自己的语言重新组织表述。直接复制原题容易导致论文重复率较高。

\item \textbf{删除冗余内容。} 题目中的故事背景、案例介绍、新闻材料等一般无需全部保留，只保留与建模相关的信息即可。

\item \textbf{保持原题意思不变。} 不要随意增加题目没有给出的条件，也不要提前加入模型假设或个人理解。

\item \textbf{按题目顺序整理。} 将问题一、问题二、问题三依次列出，使后续"问题分析"和"模型建立"保持一致。

\item \textbf{控制篇幅。} 一般控制在\textbf{半页至一页}即可，不必花费过多篇幅。
\end{enumerate}

\textbf{问题重述的目的，是让读者快速了解题目，而不是展示建模能力，因此做到"准确、简洁、完整"即可。}

\subsection{问题提出}

%\paragraph{问题一：} 不会缩进 

\textbf{问题一：}

\textbf{问题二：}

\textbf{问题三：}
